\section{Synthetic Evaluation Protocol}
\label{sec:setup}

\paragraph{Task generator.}
The benchmark is intentionally small and fully controlled. Each example is a discrete sequence generated by the local scaffold in the repository's synthetic data module. Tokens are sampled uniformly, then one or two dependency groups may be injected. A \emph{copy} group forces a span to share the same token. A \emph{progression} group enforces an arithmetic-style progression modulo the vocabulary size. These two patterns are simple, but they make cross-token dependency structure explicit and measurable.

\paragraph{Default setting.}
The main experiments use 192 sequences of length 24, vocabulary size 10, and four diffusion steps. Dependency groups appear with probability 0.8 and have span lengths between 2 and 5 tokens. The long-context stress test increases the sequence length to 40, the dependency probability to 0.9, and the span length range to 3--7.

\paragraph{Metrics.}
We report overall token accuracy, dependency-token accuracy, dependency-group accuracy, non-dependency accuracy, and the average fraction of tokens that receive non-factorized treatment. Dependency-group accuracy is especially important because a region is counted as correct only if every token in the group is recovered correctly. This metric penalizes partial recovery of coupled spans.

\paragraph{Systems compared.}
Our main comparisons are:
\begin{itemize}
\item a fully factorized baseline;
\item strict-budget \method\ runs at 5\% and 10\% non-factorized token budgets;
\item a relaxed span-4 \method\ run with a 20\% cap;
\item matched-budget fixed surrogates that allow only \coupled\ or only \micro;
\item and a random matched-budget router.
\end{itemize}

\paragraph{Budget discipline.}
Every comparison in Section~\ref{sec:results} is budget matched. The method is not allowed to claim gains by simply spending unbounded extra compute. Under the strict 10\% setting, the actual non-factorized token fraction is only 2.93\%; under the relaxed span-4 setting it is 10.89\%.
For the relaxed-budget ablations, the \coupled-only, \micro-only, random, and full \method\ systems all operate under the same top-level budget cap, which makes the comparison about allocation policy rather than raw compute.

\begin{table}[t]
\centering
\caption{Main synthetic runs. The two-mode 10\% system already improves over fully factorized decoding, and the relaxed span-4 \method\ run improves further while still using non-factorized handling on a minority of tokens.}
\label{tab:main}
\input{\tableMainResultsPath}
\end{table}

\paragraph{Scope.}
This protocol is intentionally narrower than a real diffusion language modeling benchmark. There is no pretrained backbone, no human evaluation, and no external latency measurement. The benefit is that the routing signal, the compute budget, and the dependency-heavy regions are all directly inspectable.
